% Portada del caso transversal: un edificio de tres plantas con acelerómetros y tres registros de aceleración.
\documentclass[tikz,border=0pt]{standalone}
\usepackage{resonancia}
\begin{document}
\begin{tikzpicture}
  \clip[rounded corners=6pt] (0,0) rectangle (12,3.6);
  \shade[top color=rAzul!22, bottom color=white] (0,0) rectangle (12,3.6);
  % edificio
  \draw[trazo] (0.5,0.35) -- (3.1,0.35);
  \fill[apoyo] (0.5,0.2) rectangle (3.1,0.35);
  \foreach \y in {0.95,1.7,2.45} {
    \draw[fill=rGrisClaro, trazo] (0.85,\y) rectangle (2.75,{\y+0.32});
    \fill[rNaranja] (1.8,{\y+0.32}) -- (1.7,{\y+0.52}) -- (1.9,{\y+0.52}) -- cycle;
  }
  \draw[trazo] (0.95,0.35) -- (0.95,0.95);
  \draw[trazo] (2.65,0.35) -- (2.65,0.95);
  % tres registros
  \foreach \i/\y/\f/\col in {1/2.7/3.0/rAzul, 2/1.75/2.4/rVerde, 3/0.8/1.8/rNaranja} {
    \draw[eje] (4.2,\y) -- (11.6,\y);
    \draw[color=\col, line width=0.8pt, smooth] plot[domain=4.3:11.4, samples=260] (\x, {\y + 0.28*sin(\f*600*(\x-4.3)) * (0.6+0.4*sin(130*(\x-4.3)))});
  }
  \node[etiqueta, color=rGris, anchor=west] at (4.2,3.3) {\small A · B · C};
\end{tikzpicture}
\end{document}
